\documentclass[11pt,letterpaper]{article}

\usepackage[margin=0.75in]{geometry}
\usepackage{enumitem}
\usepackage{titlesec}
\usepackage{hyperref}
\usepackage{xcolor}
\usepackage{fontenc}
\usepackage{lmodern}

% Colors
\definecolor{linkblue}{HTML}{1a0dab}

% Hyperref setup
\hypersetup{
    colorlinks=true,
    linkcolor=linkblue,
    urlcolor=linkblue,
    pdfauthor={Parker Whitfill},
    pdftitle={Parker Whitfill - Resume}
}

% Section formatting
\titleformat{\section}{\large\bfseries\scshape}{}{0em}{}[\titlerule]
\titlespacing*{\section}{0pt}{14pt}{6pt}

% Remove page numbers
\pagestyle{empty}

% Compact lists
\setlist[itemize]{nosep, left=0pt, itemsep=2pt, topsep=2pt}

\begin{document}

% ===== Header =====
\begin{center}
    {\LARGE \textbf{Parker Whitfill}} \\[6pt]
    Massachusetts Institute of Technology, Department of Economics \\
    50 Memorial Dr, Cambridge, MA 02142 \\[2pt]
    \href{mailto:parkerwhitfill@gmail.com}{parkerwhitfill@gmail.com}
    \quad $\vert$ \quad
    \href{mailto:whitfill@mit.edu}{whitfill@mit.edu}
    \quad $\vert$ \quad
    \href{https://github.com/parkerwhitfill}{github.com/parkerwhitfill}
    \quad $\vert$ \quad
    U.S.\ Citizen
\end{center}

% ===== Summary =====
\section{Summary}
Economics PhD researcher with deep expertise in AI/ML systems, benchmarking, and computational methods.
Published in \textit{Nature} on large-scale AI evaluation.
Research collaborator with Epoch AI and METR on AI forecasting, compute economics, and algorithmic progress.
Experienced in building ML pipelines for OCR, developing AI benchmarks, and full-stack software development.
Proficient in Python, R, and modern ML frameworks.

% ===== Education =====
\section{Education}

\textbf{Ph.D.\ in Economics} \hfill 2023 -- Present \\
Massachusetts Institute of Technology \hfill Cambridge, MA
\begin{itemize}
    \item Research focus: AI capabilities evaluation, benchmark design, compute economics, algorithmic progress
    \item Coursework in machine learning, econometrics, causal inference, and computational methods
\end{itemize}

\vspace{4pt}

\textbf{B.A.\ in Economics} \hfill 2017 -- 2021 \\
University of Chicago \hfill Chicago, IL

% ===== Experience =====
\section{Research \& Professional Experience}

\textbf{Future of Life Institute Fellow -- Software Engineer} \hfill 2025 \\
Future of Life Institute \hfill Cambridge, MA
\begin{itemize}
    \item Designed and built a cooperative technology platform enabling users to coordinate messaging and collective action, from architecture through deployment
    \item Developed full-stack application using Python and JavaScript, integrating APIs for real-time communication and user coordination workflows
    \item Implemented backend services, database schema design, and frontend interfaces over a 3-month intensive fellowship focused on AI governance and safety tooling
\end{itemize}

\vspace{4pt}

\textbf{Research Assistant} \hfill 2021 -- 2023 \\
University of Chicago \hfill Chicago, IL
\begin{itemize}
    \item Designed and implemented machine learning pipelines for optical character recognition (OCR) of historical documents, automating extraction of structured data from large archival corpora
    \item Built text preprocessing, image segmentation, and classification models using Python (PyTorch, OpenCV, Tesseract) to digitize and parse historical economic records
    \item Developed data ingestion and cleaning workflows for large-scale historical datasets, enabling quantitative analysis of inequality and institutional change
    \item Conducted empirical research on economic history, contributing to a publication on the causal effects of WWII on inequality and the social contract in Britain
    \item Managed end-to-end research pipelines: data collection, feature engineering, statistical modeling, and visualization
\end{itemize}

% ===== Publications =====
\section{Publications \& Research}

\textbf{A benchmark of expert-level academic questions to assess AI capabilities} \\
Coauthor; consortium paper with $>$1{,}000 authors \\
\textit{Nature}, 649:1139--1146, 2026. \href{https://www.nature.com/articles/s41586-025-09962-4}{doi:10.1038/s41586-025-09962-4}
\begin{itemize}
    \item Co-developed Humanity's Last Exam (HLE), a 3{,}000-question multi-modal benchmark evaluating frontier AI systems across dozens of academic disciplines
    \item Contributed to benchmark design, question curation, and evaluation methodology for assessing AI capabilities at expert level
\end{itemize}

\vspace{4pt}

\textbf{Do Benchmarks Overstate AI Capabilities? Evidence from SWE-Bench Verified} \\
Parker Whitfill, Cheryl Wu, Nate Rush, Joel Becker \hfill \textit{Working Paper}
\begin{itemize}
    \item Investigating measurement validity of widely-used software engineering AI benchmarks
\end{itemize}

\vspace{4pt}

\textbf{Will Compute Bottlenecks Prevent an Intelligence Explosion?} \\
Parker Whitfill and Cheryl Wu. \href{https://arxiv.org/abs/2507.23181}{arXiv:2507.23181}
\begin{itemize}
    \item Modeled compute supply-demand dynamics for AI scaling; estimated substitutability between compute and algorithmic efficiency
\end{itemize}

\vspace{4pt}

\textbf{The Software Intelligence Explosion Debate Needs Experiments} \\
Anson Ho (Epoch AI), Parker Whitfill. \href{https://epoch.ai/gradient-updates/the-software-intelligence-explosion-debate-needs-experiments}{Epoch AI}
\begin{itemize}
    \item Collaboration with Epoch AI analyzing feasibility of recursive AI self-improvement; proposed experimental frameworks for measuring AI R\&D productivity
\end{itemize}

\vspace{4pt}

\textbf{Forecasting AI Time Horizon Under Compute Slowdowns} \\
Parker Whitfill, Ben Snodin, Joel Becker (METR). \href{https://arxiv.org/abs/2511.19492}{arXiv:2511.19492}
\begin{itemize}
    \item Collaboration with METR; built quantitative forecasting models for AI development timelines under hardware constraint scenarios
\end{itemize}

\vspace{4pt}

\textbf{Note on Selection Bias in Observational Estimates of Algorithmic Progress} \\
Parker Whitfill. \href{https://arxiv.org/abs/2508.11033}{arXiv:2508.11033}

\vspace{4pt}

\textbf{Beyond Ordinal Preferences: Why Alignment Needs Cardinal Human Feedback} \\
Parker Whitfill and Stewart Slocum. \href{https://arxiv.org/abs/2508.08486}{arXiv:2508.08486}
\begin{itemize}
    \item Analyzed limitations of RLHF preference models; proposed cardinal feedback mechanisms for AI alignment
\end{itemize}

\vspace{4pt}

\textbf{The Second World War, Inequality and the Social Contract in England} \\
Leander Heldring, James Robinson, Parker Whitfill (equal contribution) \\
\textit{Economica}, Volume 89, Issue S1, Pages S137--S159

% ===== Technical Skills =====
\section{Technical Skills}

\textbf{Languages:} Python, R, SQL, \LaTeX, Bash, JavaScript/HTML/CSS \\
\textbf{ML/AI:} PyTorch, scikit-learn, OpenCV, Tesseract OCR, Hugging Face Transformers, LLM evaluation \\
\textbf{Data \& Compute:} Pandas, NumPy, Stata, large-scale data processing, statistical modeling \\
\textbf{Tools:} Git, Linux/Unix, cloud computing, Jupyter, VS Code \\
\textbf{Methods:} Machine learning, deep learning, computer vision (OCR), NLP, causal inference, econometrics, benchmark design \& evaluation, time-series forecasting

% ===== Honors =====
\section{Honors \& Fellowships}
\begin{itemize}
    \item Emergent Ventures Grantee (2025)
    \item Stripe Economics of AI Fellowship (2025)
    \item Future of Life Institute Fellowship (2025)
    \item 2017 National Debate Champion (TOC LD)
\end{itemize}

\end{document}
